%% sections/06-conclusion.tex

\section{Conclusion}
\label{sec:conclusion}

We compared two resource exhaustion mechanics --- Campaign 4's Supply Die
(degrading from d12 to exhausted over 7 sessions) and Campaign 6's Unit Cohesion
Die (degrading d12 to d8 as unit members were lost) --- against five campaigns with
no degrading-resource mechanics. Both resource campaigns produced monotone-rising
Character Agency trajectories, reaching the rubric ceiling (CA 10) by session 3--4
and maintaining it through the finale. Non-resource campaigns showed late-rising
trajectories, typically reaching CA 10 only at session 7.

The Supply Die produced a +0.6 CA-point advantage over light-pressure baselines
at sessions 1--6, attributable to earlier arrival at the CA 10 ceiling. The Unit
Cohesion Die produced a comparable gate-score advantage, with an additional +3.0
panel-session advantage in sessions 5--7 (three consecutive sessions rated
\textit{character-primary} by the Tactician lens vs.\ \textit{character-inflected
tactical} for the Supply Die). The additional advantage of the Cohesion Die is
explained by its tracking of named people rather than materials: the party preserved
the Cohesion resource because they cared about the people it represented, not because
preservation was tactically optimal.

Our central argument is that resource exhaustion amplifies character-driven
decision-making in long-form campaigns when two conditions are met: the resource
is \textit{legible} (the party can see its depletion) and \textit{emotionally
architectured} (the resource represents something the party has reason to care about
beyond tactical survival). When both conditions hold, resource pressure does not
crowd out character-driven engagement; it channels character values into resource
decisions, making those decisions the primary medium through which the party expresses
who they are.

The finding inverts the conventional assumption that tactical pressure suppresses
narrative engagement. In short-form games, this assumption may hold. In long-form
campaigns (7+ sessions), resource depletion accumulates a history, and that history
becomes emotionally freighted: the party's decisions are made in the context of
everything that has already been spent and everyone who has already been lost.
It is precisely when the resources are most depleted --- when the die is at d4,
when the unit is at d8 --- that the party's character-driven engagement is highest,
because those are the sessions when every decision is simultaneously a tactical
choice and a statement about what the party is willing to do.

\subsection{Future Work}

Three directions are most pressing.

\textbf{Controlled archetype comparison.} The critical next step is a controlled
experiment: two campaigns of the same archetype, one with and one without a degrading-resource
mechanic, run in the same system. The Marathon pipeline supports this comparison
directly --- the same adventure module can be run with and without the Supply Die
or Cohesion Die active. This comparison would resolve the archetype confound that
is the primary limitation of the current analysis.

\textbf{Human player replication.} AI-simulated sessions may not capture human
player behavior accurately, particularly in the grief-paragraph channel (whether
players engage character grief through resource decisions in the same pattern that
the simulation produces). Running the same campaigns with human players and comparing
CA trajectories would establish whether the amplification effect is robust to player
type.

\textbf{Cross-system generalization.} The mechanics studied here are D\&D 5e-specific
(hit dice, spell slots, action economy). Replicating the study in a system like
Ironsworn (where supply and spirit tracks are explicitly resource mechanics) or
Blades in the Dark (where stress and trauma are the primary attrition system) would
test whether the amplification hypothesis is system-independent or specific to D\&D's
resource architecture.

\subsection{Design Recommendations}

For practitioners designing long-form narrative campaigns with resource mechanics:
implement a visible degrading die (not invisible HP math) that tracks something the
party will care about losing --- preferably named people rather than abstract
materials. Do not allow full resource recovery between sessions; the cumulative
depletion is the mechanism, not the instantaneous pressure. Allow active preservation
to be a meaningful character choice by designing encounters where conservation is
possible, if not tactically optimal. And commit to the 5--7 session form: the
amplification effect accumulates over sessions, and campaigns shorter than 5 sessions
may not allow enough depletion history for the emotional architecture to become load-bearing.

Resource exhaustion, designed this way, is not an obstacle to narrative engagement.
It is its engine.
